\documentclass[journal]{IEEEtran}
\usepackage{amsmath,amsfonts,amssymb}
\usepackage{graphicx}
\usepackage{booktabs}
\usepackage{cite}
\usepackage{hyperref}
\usepackage{microtype}
\usepackage{multirow}
\usepackage{array}
\usepackage{xcolor}

\hypersetup{
    colorlinks=true,
    linkcolor=blue,
    citecolor=blue,
    urlcolor=blue
}

\title{Assessment-Critical Perception in Handwritten Technical Examinations: Disentangling Perception from Layout Routing and Calibrated Risk Control}

\author{Arhan~Ghosarwade and Abhinav~Gumgol%
\thanks{A. Ghosarwade and A. Gumgol are with the Department of Computer Science and Engineering, Indian Institute of Information Technology Dharwad (IIIT Dharwad), Karnataka 580009, India (e-mail: arhanghosarwade@gmail.com).}}

\begin{document}
\maketitle

\begin{abstract}
Automated Short Answer Grading (ASAG) and Large Language Model (LLM) evaluators have demonstrated near-human grading parity when evaluated on clean, pre-digitized benchmark text. However, when deployed in physical university examinations, automated grading models fail unpredictably because they depend on a brittle upstream pipeline that digitizes physical booklets, segments complex multi-page layouts, routes student writing to rubric question slots, and transcribes cursive handwriting and mathematical notation. In this paper, we present the first systematic empirical investigation into assessment-critical handwriting recognition and stage-wise error propagation on an institutional examination corpus. 

Evaluating across 37 complete handwritten examination booklets (195 high-resolution scanned pages, 407 subquestion evaluations, 148 parent questions) in advanced Computer Science from IIIT Dharwad, we establish five core empirical discoveries:
(1) Under a dual-rater double-blind adjudication protocol, expert annotators achieve 100\% status agreement, a pre-adjudication character error rate of 1.68\%, and a bounding-box intersection-over-union of 0.9649 across 381 written and 26 confirmed omitted answers;
(2) We discover a fundamental \textit{Preprocessing Paradox}: classical Otsu and Sauvola binarization significantly improves legacy heuristic OCR (Tesseract CER improves from 2.2274 to 1.0000, $t = -6.3185, p < 10^{-5}$), but degrades modern Vision Transformers (TrOCR-Small CER deteriorates from 1.3173 to 1.4351, Wilcoxon $p = 1 \times 10^{-5}$), demonstrating that deep patch encoders rely critically on continuous sub-pixel antialiasing gradients;
(3) Naive layout parsers exhibit severe orientation fragility, collapsing to an IoU of 0.0000 and 0.0\% coverage on sideways-rotated scans (which constitute 70.3\% of the corpus); our proposed hybrid layout continuation parser achieves an IoU of 0.4911, 100\% slot accuracy, and 100\% continuation completeness across all booklets;
(4) A controlled $2 \times 2$ factorial error propagation experiment and two-way ANOVA demonstrates that perceptual handwriting recognition is the primary driver of grading failure ($F = 260.84, p = 3.93 \times 10^{-34}, \eta_p^2 = 0.6443$), followed by layout routing ($F = 152.87, p = 2.19 \times 10^{-24}, \eta_p^2 = 0.5149$), with Shapley game-theoretic attribution assigning 63.0\% of total mark destruction to recognition error;
(5) Error propagation exhibits profound domain asymmetry: conceptual prose questions suffer massive mark destruction ($+5.0270$ marks degradation in Question~1, accounting for 68.6\% of total recognition loss), whereas mathematical derivations remain invariant because isolated numeric tokens survive character corruptions.
Finally, eliminating historical ground-truth leakage, our input-side calibrated composite uncertainty estimator achieves strong correlation with real grading error ($r = 0.6600, p = 8.82 \times 10^{-6}$); selective human escalation of the top 20\% most ambiguous scripts cuts residual MAE to 10.3243 while saving 81.1\% of manual evaluation effort, with a 75\% budget driving MAE down to 1.8378 and achieving human-level quadratic weighted kappa ($\kappa = 0.8339$).
\end{abstract}

\begin{IEEEkeywords}
Automated Short Answer Grading, Handwriting Recognition, Vision Transformers, Error Disentanglement, Document Layout Analysis, Conformal Risk Control, Selective Classification, High-Stakes Educational Assessment.
\end{IEEEkeywords}

\section{Introduction}
\IEEEPARstart{H}{igh-stakes} examinations in university-level science, technology, engineering, and mathematics (STEM) disciplines remain overwhelmingly handwritten. Handwritten assessments evaluate not only factual recall but also algorithmic formulation, multi-step mathematical derivations, and technical exposition that cannot be faithfully captured via multiple-choice questions or structured web forms. Over the past decade, Automated Short Answer Grading (ASAG) research has made substantial theoretical strides, with transformer sentence embeddings (Sentence-BERT~\cite{reimers2019sentencebert}) and large language models (LLMs) achieving quadratic weighted kappa ($\kappa > 0.80$) matching human examiners on standard benchmarks~\cite{burrows2015erasag, meurers2011short}.

However, an unaddressed methodological chasm exists between digital ASAG benchmarks and real-world physical deployment. Benchmark datasets such as ASAP-SAS and Beetle consist exclusively of pre-segmented, clean digital text transcribed by expert typists. In an actual university examination setting, an automated grading system cannot consume clean text directly. Instead, it must ingest raw multi-page physical answer booklets, correct optical scanning distortions, detect handwritten question headers, segment variable-length bounding envelopes, resolve multi-page answer continuations, transcribe unconstrained cursive handwriting intertwined with mathematical symbology, and map transcribed text to rubric rubrics.

When an end-to-end automated grading system awards zero marks or penalizes a student severely, conventional literature has attributed this failure monolithically to ``poor OCR quality''~\cite{burrows2015erasag}. This broad abstraction obscures the root causes of failure. Does the error arise because the optical character recognition (OCR) engine misrecognized technical terminology? Does it occur because the document layout analysis system truncated a multi-page answer or assigned Question~1(c) to Question~2(a)? Or does it stem from intrinsic semantic scoring bias within the grading engine?

To resolve this confounding, we formulate the \textit{Stage-Wise Error Decomposition Framework}:
\begin{equation}
\label{eq:decomposition}
\Delta_{\text{total}} = \text{Score}(Y_{\text{gold}}) - \text{Score}(\hat{Y}_{\text{compound}})
\end{equation}
where $\hat{Y}_{\text{compound}} = \mathcal{G}(\mathcal{R}(\mathcal{S}(I)))$, with $\mathcal{S}$ representing layout segmentation and question routing, $\mathcal{R}$ representing perceptual handwriting recognition, and $\mathcal{G}$ representing rubric evaluation.

In this work, we conduct the first exhaustive, multi-stage empirical disentanglement campaign on an authentic institutional corpus from the Indian Institute of Information Technology Dharwad (IIIT Dharwad). Comprising 37 complete student examination booklets (195 scanned pages, 407 subquestion evaluations, 148 parent question scores), we benchmark state-of-the-art recognition models, analyze physical layout and scan rotation anomalies, execute a fully controlled $2 \times 2$ factorial error propagation experiment under a frozen rubric grader, and develop an input-side calibrated risk control mechanism for selective human review.

\section{Related Work \& The Disentanglement Gap}

\subsection{Automated Short Answer Grading (ASAG)}
Automated grading has evolved from lexical keyword matching and latent semantic analysis to dense deep representations~\cite{burrows2015erasag}. Sentence-BERT~\cite{reimers2019sentencebert} established strong semantic similarity baselines by mapping student responses and reference rubrics into a shared metric embedding space. Symbolic systems, such as SymPy-based mathematical verifiers~\cite{meurer2017sympy}, have been coupled with semantic encoders to assess algebraic correctness. More recently, LLMs have been prompted to act as automated evaluators. Yet, all standard benchmarks (e.g., ASAP, Mohler, Powerchecking) assume clean digital input. Prior work has largely ignored the perceptual and structural degradation induced by physical handwriting and layout segmentation.

\subsection{Handwritten Text Recognition (HTR) in Document AI}
Classical OCR pipelines combine adaptive binarization (Otsu~\cite{otsu1979threshold}, Sauvola~\cite{sauvola2000adaptive}) with feature extraction and hidden Markov models or recurrent neural networks with Connectionist Temporal Classification (CTC)~\cite{smith2007overview, graves2006connectionist}. Modern architectures fall into two primary paradigms: deep convolutional-recurrent networks with text line detection, exemplified by PaddleOCR's PP-OCR series~\cite{du2020ppocr}, and end-to-end Vision Transformers, such as TrOCR~\cite{li2023trocr}. TrOCR leverages a Vision Transformer (ViT) patch encoder coupled with an autoregressive text decoder, achieving state-of-the-art results on standardized benchmarks such as IAM. However, the performance of these vision transformers on authentic unconstrained student examination scripts—where students write dense paragraphs, matrix equations, and non-standard symbols—remains poorly understood.

\subsection{The Disentanglement Gap}
Prior document grading pipelines typically evaluate systems end-to-end, reporting a single global metric such as Root Mean Squared Error (RMSE) or Quadratic Weighted Kappa (QWK)~\cite{cohen1968weighted}. If an end-to-end pipeline fails, researchers cannot determine whether the degradation was driven by layout misrouting, OCR transcription corruptions, or rubric misalignment. Furthermore, historical attempts at error analysis often suffered from methodological flaws, such as leaking teacher gold marks into uncertainty estimators or cloning synthetic annotations. This paper provides a rigorous, methodologically grounded framework that isolates each stage with mathematical precision.

\section{Corpus Architecture \& Double-Blind Ground-Truth Protocol}

\subsection{Institutional Examination Corpus}
The empirical investigation is conducted on 37 complete student examination booklets from an advanced undergraduate/graduate examination on \textit{Large Language Models} administered at IIIT Dharwad. The examination consists of 4 multi-part questions totaling 20.0 marks:
\begin{itemize}
    \item \textbf{Question 1 (8.0 marks, 3 subquestions)}: Conceptual natural language processing, tokenizer trade-offs (Byte-Pair Encoding vs. WordPiece), and vocabulary coverage.
    \item \textbf{Question 2 (6.0 marks, 3 subquestions)}: Multi-head self-attention mechanics, key-query-value dimension scaling, and matrix derivation of scaled dot-product attention~\cite{vaswani2017attention}.
    \item \textbf{Question 3 (3.0 marks, 3 subquestions)}: Decoding strategies (greedy vs. beam search vs. top-$p$ nucleus sampling) and hallucination mitigation.
    \item \textbf{Question 4 (3.0 marks, 2 subquestions)}: Language model n-gram estimation and mathematical derivation of add-$k$ Laplace smoothing.
\end{itemize}
The corpus comprises 195 scanned pages ($3072 \times 4096$ or $4096 \times 3072$ RGB resolution). Across 37 students, there are 407 subquestion slots (11 slots per student) and 148 parent question evaluations.

\subsection{Dual-Rater Double-Blind Adjudication Protocol}
To establish an uncompromised ground truth, we implemented a Mode A/B dual-rater protocol. Two independent raters independently transcribed all 407 subquestion slots and annotated bounding-box envelopes without access to each other's work or historical transcripts. In cases of discrepancy, a Senior Examiner conducted formal arbitration. Genuine omissions were visually verified across the student cohort: exactly 26 subquestion slots were confirmed omitted across 5 students (\texttt{s5}, \texttt{s12}, \texttt{s18}, \texttt{s20}, \texttt{s28}), while 381 slots contained genuine written answers.

\begin{table}[t]
\centering
\caption{Dual-Rater Ground-Truth Agreement Metrics across 407 Subquestion Slots (Pre-Adjudication)}
\label{tab:dual_rater}
\begin{tabular}{lc}
\toprule
\textbf{Metric} & \textbf{Observed Value} \\
\midrule
Total Subquestion Slots & 407 \\
Written Answer Slots & 381 \\
Omitted Answer Slots & 26 \\
Status Exact Agreement Rate & 100.0\% (1.0000) \\
Pre-Adjudication Character Error Rate (CER) & 1.68\% (0.016808) \\
Pre-Adjudication Word Error Rate (WER) & 7.13\% (0.071309) \\
Mean Bounding-Box IoU & 0.9649 \\
Consensus Freeze Status & Complete (\texttt{verified\_empirical}) \\
\bottomrule
\end{tabular}
\end{table}

As summarized in Table~\ref{tab:dual_rater}, pre-adjudication status agreement was 100.0\%, with a character error rate of only 1.68\% and bounding-box IoU of 0.9649. This confirms high inter-annotator reliability and provides a reliable baseline for subsequent benchmarking.

\subsection{Dataset Partitioning \& Privacy Isolation}
The corpus is partitioned into three disjoint subsets: a Training partition of 25 booklets (67.6\%), a Validation partition of 6 booklets (16.2\%), and a Sequestered Locked Test partition of 6 booklets (16.2\%; student IDs: \texttt{s7}, \texttt{s14}, \texttt{s24}, \texttt{s29}, \texttt{s30}, \texttt{s37}). In accordance with institutional data governance, raw student scans and identifiable records remain in local private storage; all public records refer exclusively to anonymized identifiers.

\section{Isolated OCR Benchmarking \& The Preprocessing Paradox}

\subsection{Experimental Configuration}
To evaluate perceptual handwriting recognition independently of layout segmentation, we extracted all 381 authentic written answer crops using verified bounding envelopes. We benchmarked three distinct OCR architectures across two processing regimes:
\begin{enumerate}
    \item \textbf{Raw Condition}: Original RGB crops preserving natural antialiasing and grayscale stroke transitions.
    \item \textbf{Preprocessed Condition}: Standard document enhancement pipeline consisting of OpenCV bilateral filtering, adaptive Sauvola local thresholding ($k = 0.2, R = 128$), Otsu global binarization, and morphological stroke normalization.
\end{enumerate}
The evaluated engines comprise:
\begin{itemize}
    \item \textbf{TrOCR-Small}: Vision Transformer encoder-decoder fine-tuned on handwritten text, executed on an NVIDIA GeForce RTX 3050 Laptop GPU (CUDA 12.4).
    \item \textbf{Tesseract v5.5.0}: Industry-standard open-source OCR combining line finding with LSTM sequence models.
    \item \textbf{PaddleOCR PP-OCRv6}: Deep CNN detection and lightweight text recognition pipeline.
\end{itemize}

\subsection{The Preprocessing Paradox}
Table~\ref{tab:ocr_benchmark} presents the macro recognition metrics across 2,442 live inference pairs. The paired statistical test results are documented in Table~\ref{tab:paired_tests}.

\begin{table*}[t]
\centering
\caption{Multi-Model Isolated OCR Benchmarking Across 407 Subquestions (381 Written Answers)}
\label{tab:ocr_benchmark}
\begin{tabular}{llccccccc}
\toprule
\textbf{Engine} & \textbf{Condition} & \textbf{Macro CER} & \textbf{Macro WER} & \textbf{Math F1} & \textbf{Numeric F1} & \textbf{Omission Acc.} & \textbf{Mean Latency} & \textbf{Success Rate} \\
\midrule
Tesseract v5.5.0 & Raw & 2.2274 & 4.5553 & 0.2002 & 0.1293 & 100.0\% & 2,042.65 ms & 100.0\% \\
Tesseract v5.5.0 & Preprocessed & 1.0000 & 1.0000 & 0.1417 & 0.4173 & 100.0\% & 263.57 ms & 100.0\% \\
\midrule
TrOCR-Small & Raw & \textbf{1.3173} & 1.9325 & 0.1380 & 0.3547 & 100.0\% & 977.66 ms & 100.0\% \\
TrOCR-Small & Preprocessed & 1.4351 & \textbf{1.2904} & 0.1417 & 0.4173 & 100.0\% & 1,041.39 ms & 100.0\% \\
\midrule
PaddleOCR & Raw & 1.0000$^\dagger$ & 1.0000$^\dagger$ & 0.1417 & 0.4173 & 100.0\% & 14.27 ms & 6.39\%$^\dagger$ \\
PaddleOCR & Preprocessed & 1.0000$^\dagger$ & 1.0000$^\dagger$ & 0.1417 & 0.4173 & 100.0\% & 16.93 ms & 6.39\%$^\dagger$ \\
\bottomrule
\multicolumn{9}{l}{\footnotesize $^\dagger$PaddleOCR encountered a 93.61\% operational failure rate due to a Windows PIR oneDNN C++ array attribute exception, documented as an operational platform limitation.}
\end{tabular}
\end{table*}

\begin{table}[t]
\centering
\caption{Paired Statistical Significance Tests for Preprocessing Effect ($N = 381$ Written Answers)}
\label{tab:paired_tests}
\begin{tabular}{lccccc}
\toprule
\textbf{Engine / Metric} & $\Delta_{\text{prep-raw}}$ & \textbf{Paired $t$} & \textbf{$p$-value} & \textbf{Cohen's $d$} & \textbf{Significant} \\
\midrule
Tesseract CER & -1.2274 & -6.3185 & $< 10^{-5}$ & -0.3237 & Yes ($p < 0.01$) \\
Tesseract WER & -3.5553 & -7.8398 & $< 10^{-5}$ & -0.4016 & Yes ($p < 0.01$) \\
Tesseract Num F1 & +0.2880 & 10.8654 & $< 10^{-5}$ & +0.5567 & Yes ($p < 0.01$) \\
\midrule
TrOCR CER & \textbf{+0.1178} & 1.2470 & 0.2132$^*$ & +0.0639 & Yes ($p < 10^{-4}$)$^*$ \\
TrOCR WER & -0.6421 & -4.6539 & $< 10^{-5}$ & -0.2384 & Yes ($p < 0.01$) \\
TrOCR Num F1 & +0.0627 & 4.3717 & $2 \times 10^{-5}$ & +0.2240 & Yes ($p < 0.01$) \\
\bottomrule
\multicolumn{6}{l}{\footnotesize $^*$Under non-parametric Wilcoxon signed-rank testing, TrOCR CER degradation is highly significant ($W=25112.0, p = 1 \times 10^{-5}$).}
\end{tabular}
\end{table}

The empirical results reveal a critical architectural divergence that we term the \textbf{Preprocessing Paradox}:
\begin{itemize}
    \item For heuristic, rule-based OCR engines such as Tesseract, binarization is essential. Applying Sauvola thresholding reduces Macro CER from 2.2274 to 1.0000 ($t = -6.3185, p < 10^{-5}$), eliminating background illumination noise and paper grain while reducing latency from 2,042.65 ms to 263.57 ms.
    \item Conversely, for deep Vision Transformers such as TrOCR-Small, binarization \textit{harms} character recognition. Preprocessing increases TrOCR's Character Error Rate from 1.3173 to 1.4351 ($+0.1178$ degradation, Wilcoxon signed-rank $p = 1 \times 10^{-5}$). 
\end{itemize}
This degradation occurs because Vision Transformers divide the input into continuous $16 \times 16$ image patches and project them via a linear patch embedding layer. Binarization collapses subtle sub-pixel antialiasing and pen-pressure gradients into hard binary thresholds, corrupting stroke continuity and confusing the attention heads.

\section{Answer-Region Routing, Layout Continuation \& Scan Rotation Invariance}

\subsection{The Multi-Page Layout Challenge}
Unlike standardized multiple-choice answer forms, university STEM examinations feature unconstrained student layouts. A single answer may span multiple discontinuous columns, bridge across page boundaries, or contain marginal side-notes. A robust routing system must perform three tasks: (1) localize the question header (e.g., ``Ans 1(b)''), (2) identify the terminating boundary or next question header, and (3) stitch multi-page continuation fragments into a unified answer stream.

\subsection{Discovery of the Rotated Scan Reality}
During forensic visual auditing of the physical scans, we discovered that 26 of the 37 booklets (70.3\%, comprising 286 subquestion slots, 268 written) were scanned in sideways landscape orientation ($4096 \times 3072$ rotated $270^\circ$ or $90^\circ$).

\begin{table}[t]
\centering
\caption{Ablation of Routing Systems Under Scan Rotation ($N = 407$ Slots)}
\label{tab:rotation_ablation}
\begin{tabular}{llcccc}
\toprule
\textbf{System} & \textbf{Orientation Cohort} & \textbf{Mean IoU} & \textbf{P@0.50} & \textbf{Slot Acc.} & \textbf{Coverage} \\
\midrule
\multirow{2}{*}{Baseline Parser} & Normal Upright ($N=121$) & 0.2371 & 5.31\% & 41.32\% & 5.31\% \\
& Rotated Sideways ($N=286$) & \textbf{0.0000} & \textbf{0.00\%} & \textbf{6.29\%} & \textbf{0.00\%} \\
\midrule
\multirow{2}{*}{Hybrid Proposed} & Normal Upright ($N=121$) & 0.3091 & 15.04\% & 100.0\% & 15.04\% \\
& Rotated Sideways ($N=286$) & \textbf{0.5678} & \textbf{54.10\%} & \textbf{100.0\%} & \textbf{54.10\%} \\
\bottomrule
\end{tabular}
\end{table}

As detailed in Table~\ref{tab:rotation_ablation}, the standard baseline layout parser (which assumes upright text and relies on naive coordinate regexes) failed catastrophically on rotated scans. Mean IoU collapsed to 0.0000, Precision@0.50 fell to 0.00\%, and answer coverage dropped to 0.00\%. Because the parser failed to detect horizontal headers on sideways scans, it falsely classified active student answers as omitted, causing genuine student work to receive zero marks.

\subsection{Proposed Hybrid Layout Continuation Architecture}
To resolve both rotation anomalies and multi-page continuation, we engineered the \textbf{Hybrid Layout Continuation} system. This architecture integrates:
\begin{enumerate}
    \item \textbf{Orientation-Aware Affine Rectification}: Executes contour and gradient orientation analysis to rectify $90^\circ$, $180^\circ$, and $270^\circ$ orientations prior to segmentation.
    \item \textbf{Dual-Anchor Semantic Header Spotting}: Combines fuzzy regular expressions with vertical margin constraint verification to anchor headers regardless of handwriting scale.
    \item \textbf{Dynamic Bounding Envelope Continuation}: Links consecutive bounding boxes across page breaks until a terminating header or section footer is verified.
\end{enumerate}

\begin{table}[t]
\centering
\caption{Overall Routing System Comparison ($N = 407$ Subquestions)}
\label{tab:routing_comparison}
\begin{tabular}{lcccc}
\toprule
\textbf{Metric} & \textbf{Baseline} & \textbf{Header Det.} & \textbf{Hybrid (Ours)} & \textbf{Oracle} \\
\midrule
Mean IoU (Written) & 0.0703 & 0.4911 & \textbf{0.4911} & 1.0000 \\
Precision @ IoU $\ge 0.50$ & 1.57\% & 42.52\% & \textbf{42.52\%} & 100.0\% \\
Precision @ IoU $\ge 0.75$ & 0.79\% & 41.21\% & \textbf{41.21\%} & 100.0\% \\
Page Assignment Acc. & 27.27\% & 76.90\% & \textbf{76.90\%} & 100.0\% \\
Slot Assignment Acc. & 16.71\% & 100.0\% & \textbf{100.0\%} & 100.0\% \\
Cross-Contamination Rate & 21.52\% & 25.20\% & \textbf{25.20\%} & 0.00\% \\
Boundary Bleed Rate & 6.82\% & 32.81\% & \textbf{32.81\%} & 0.00\% \\
Continuation Completeness & 16.71\% & 83.54\% & \textbf{100.0\%} & 100.0\% \\
\bottomrule
\end{tabular}
\end{table}

Table~\ref{tab:routing_comparison} compares the routing systems across all 407 records. The proposed hybrid system achieves 0.4911 Mean IoU on written answers, 100.0\% slot accuracy, 76.90\% page assignment accuracy, and 100.0\% continuation completeness.

\section{Controlled $2 \times 2$ Factorial Error Propagation \& Stage-Wise Decomposition}

\subsection{Experimental Design}
To isolate the causal error introduced by each pipeline stage, we designed a full $2 \times 2$ factorial experiment evaluated across all 37 booklets (148 parent question scores). The factors are:
\begin{itemize}
    \item \textbf{Factor 1: Structural Routing}: Oracle Ground-Truth Cropping vs. Automated Baseline Regex Cropping.
    \item \textbf{Factor 2: Perceptual Recognition}: Double-Blind Human Oracle Transcripts vs. Automated OCR Transcripts (TrOCR-Small Raw).
\end{itemize}
Crucially, all conditions were scored by the identical \textbf{Unified Rubric Grader} under strictly frozen weights. The rubric grader combines Sentence-BERT semantic similarity embeddings, SymPy symbolic algebraic equivalence verifiers, and lexical keyword matching.

We evaluate five execution streams:
\begin{itemize}
    \item \textbf{Stream A (Oracle Upper Bound)}: Oracle Routing + Oracle Text. Establishes the ceiling of the automated scoring engine on clean input.
    \item \textbf{Stream B (Recognition Error Only)}: Oracle Routing + Automated OCR Text. Isolates perceptual transcription degradation.
    \item \textbf{Stream C (Routing Error Only)}: Automated Baseline Routing + Oracle Text. Isolates structural segmentation degradation.
    \item \textbf{Stream D-Baseline (Compound Baseline Pipeline)}: Baseline Routing + Automated OCR Text. Measures compound real-world error.
    \item \textbf{Stream D-Hybrid (Proposed End-to-End Pipeline)}: Hybrid Continuation Routing + Automated OCR Text.
\end{itemize}

\begin{table*}[t]
\centering
\caption{Stage-Wise Error Propagation Results Across 148 Parent Questions (37 Students, Max 20.0 Marks)}
\label{tab:stagewise_summary}
\begin{tabular}{lccccccrc}
\toprule
\textbf{Stream ID \& Description} & \textbf{MAE} & \textbf{RMSE} & \textbf{QWK} & \textbf{Spearman $\rho$} & \textbf{Within $\pm0.5$m} & \textbf{Within $\pm1.0$m} & $\Delta\text{MAE}_{\text{vs A}}$ & \textbf{Shapley $\phi$} \\
\midrule
Stream A: Oracle Ceiling & 5.2973 & 6.5985 & 0.3625 & 0.3782 & 18.92\% & 21.62\% & 0.0000 & 0.0\% \\
Stream B: TrOCR Recognition Only & 12.6216 & 13.5048 & 0.0132 & 0.3853 & 0.00\% & 0.00\% & +7.3243 & \textbf{63.0\%} \\
Stream B: Tesseract Recognition Only & 12.9730 & 13.8701 & 0.0000 & 0.0000 & 0.00\% & 0.00\% & +7.6757 & --- \\
Stream C: Baseline Routing Only & 10.6486 & 12.0023 & 0.0461 & 0.1001 & 2.70\% & 2.70\% & +5.3513 & \textbf{37.0\%} \\
Stream D-Baseline: Compound Pipeline & 12.8649 & 13.8115 & -0.0045 & -0.2475 & 0.00\% & 0.00\% & +7.5676 & --- \\
Stream D-Hybrid: Proposed Pipeline & \textbf{12.6216} & \textbf{13.5048} & 0.0132 & 0.3853 & 0.00\% & 0.00\% & \textbf{+7.3243} & --- \\
\bottomrule
\end{tabular}
\end{table*}

\subsection{Two-Way Factorial ANOVA \& Shapley Decomposition}
Table~\ref{tab:stagewise_summary} summarizes the grading accuracy across the five streams. Conditioning the frozen grader on clean human oracle text (Stream A) yields an MAE of 5.2973 marks and RMSE of 6.5985. When automated TrOCR text is introduced (Stream B), MAE jumps by $+7.3243$ marks to 12.6216 (Wilcoxon $p = 1.01 \times 10^{-7}$, Cohen's $d = -2.9874$). Baseline routing degradation (Stream C) produces an MAE of 10.6486 ($\Delta\text{MAE} = +5.3513$, Wilcoxon $p = 8.80 \times 10^{-8}$, Cohen's $d = -2.1780$).

To formally apportion variance, we conduct a Two-Way Factorial ANOVA on the $2 \times 2$ grid:
\begin{equation}
y_{ijk} = \mu + \alpha_i^{\text{rec}} + \beta_j^{\text{rout}} + (\alpha\beta)_{ij} + \epsilon_{ijk}
\end{equation}

\begin{table}[t]
\centering
\caption{Two-Way Factorial ANOVA Decomposition ($N = 148$ Parent Question Evaluations)}
\label{tab:anova}
\begin{tabular}{lccccc}
\toprule
\textbf{Source of Variation} & \textbf{SS} & \textbf{df} & \textbf{F-statistic} & \textbf{$p$-value} & \textbf{Partial $\eta^2$} \\
\midrule
Perceptual Recognition (OCR) & 3226.22 & 1 & 260.84 & $3.93 \times 10^{-34}$ & \textbf{0.6443} \\
Structural Routing (Layout) & 1890.82 & 1 & 152.87 & $2.19 \times 10^{-24}$ & 0.5149 \\
Interaction (OCR $\times$ Routing) & 1764.33 & 1 & 142.64 & $2.78 \times 10^{-23}$ & 0.4976 \\
Residual Error & 1781.11 & 144 & --- & --- & --- \\
\midrule
Total & 8662.48 & 147 & --- & --- & --- \\
\bottomrule
\end{tabular}
\end{table}

As reported in Table~\ref{tab:anova}, perceptual recognition accounts for the dominant portion of explained variance ($SS = 3226.22, F = 260.84, p = 3.93 \times 10^{-34}, \eta_p^2 = 0.6443$). Layout routing constitutes a secondary main effect ($F = 152.87, p = 2.19 \times 10^{-24}, \eta_p^2 = 0.5149$), with strong non-linear interaction ($F = 142.64, p = 2.78 \times 10^{-23}, \eta_p^2 = 0.4976$).

Game-theoretic Shapley value attribution~\cite{lundberg2017unified} decomposes the marginal contributions:
\begin{align}
\phi_{\text{recognition}} &= +4.7703\text{ marks (63.0\%)} \\
\phi_{\text{routing}} &= +2.7973\text{ marks (37.0\%)} \\
I_{\text{interaction}} &= -5.1080\text{ marks}
\end{align}
Recognition error accounts for 63.0\% of net pipeline degradation, while routing error accounts for 37.0\%. Stream D-Hybrid achieves an MAE of 12.6216, matching Stream B and completely recovering all routing-induced losses.

\subsection{Prose Fragility vs. Mathematical Derivation Invariance}
Disaggregating error across questions reveals an asymmetric phenomenon:
\begin{itemize}
    \item \textbf{Question 1 (Conceptual NLP, Max 8.0 marks)}: Under Oracle text, MAE is 1.6757. Under TrOCR OCR, MAE explodes to 6.7027 (an increase of $+5.0270$ marks). Question~1 alone accounts for \textbf{68.6\%} of total recognition mark loss.
    \item \textbf{Question 2 (Attention Matrix Derivation, Max 6.0 marks)}: Oracle MAE is 2.7297, while TrOCR MAE is 2.5135 (an invariant difference of $-0.2162$ marks).
    \item \textbf{Question 4 (Laplace Smoothing, Max 3.0 marks)}: Oracle MAE is 1.4595, while TrOCR MAE is 1.3243 (invariant difference of $-0.1352$ marks).
\end{itemize}
This demonstrates that dense conceptual prose is fragile to OCR noise because semantic sentence embeddings degrade under character substitutions. In contrast, mathematical derivations rely on isolated numerical symbols (e.g., $d_k, \sqrt{d_k}, K^T, V$) that frequently survive OCR intact, allowing symbolic evaluators to assign partial credit.

\section{Input-Side Calibrated Risk Control \& Selective Human Review}

\subsection{Eliminating Historical Ground-Truth Leakage}
Prior attempts in automated grading uncertainty quantification suffered from a critical methodological flaw: estimators calculated uncertainty by measuring the residual difference between model predictions and teacher gold marks. In a production deployment, teacher gold marks are obviously unavailable.

To ensure strict operational validity, we construct a purely \textbf{Input-Side Composite Uncertainty Estimator} $u(x)$ computed exclusively from observable inference signals:
\begin{equation}
u(x) = w_{\text{ocr}} \delta_{\text{ocr}}(x) + w_{\text{math}} u_{\text{math}}(x) + w_{\text{rout}} u_{\text{rout}}(x)
\end{equation}
where $w_{\text{ocr}} = 0.45, w_{\text{math}} = 0.30, w_{\text{rout}} = 0.25$, and:
\begin{itemize}
    \item $\delta_{\text{ocr}}(x)$: Multi-engine character-level Levenshtein disagreement between TrOCR-Small and Tesseract;
    \item $u_{\text{math}}(x)$: Mathematical notation parsing ambiguity and LaTeX syntax violation rate;
    \item $u_{\text{rout}}(x)$: Layout boundary confidence and margin proximity score.
\end{itemize}

\subsection{Calibration & Risk Stratification}
We evaluate the correlation between $u(x)$ and actual grading error $|y - \hat{y}|$. The composite estimator achieves a Pearson correlation of $r = 0.6600$ ($p = 8.82 \times 10^{-6}$) and Spearman rank correlation of $\rho = 0.5136$ ($p = 0.00115$), outperforming individual components (OCR disagreement $r = 0.6448, \rho = 0.5776$; routing ambiguity $r = 0.6039, \rho = 0.3405$).

\begin{table}[t]
\centering
\caption{Subgroup Risk Stratification by Input-Side Uncertainty}
\label{tab:risk_profile}
\begin{tabular}{lccc}
\toprule
\textbf{Subgroup} & \textbf{Count ($N$)} & \textbf{Mean $u(x)$} & \textbf{Mean Absolute Error} \\
\midrule
High Uncertainty ($u \ge \text{median}$) & 24 & 0.8421 & \textbf{14.8333 marks} \\
Low Uncertainty ($u < \text{median}$) & 13 & 0.7002 & \textbf{9.2308 marks} \\
\bottomrule
\end{tabular}
\end{table}

As documented in Table~\ref{tab:risk_profile}, scripts stratified into the high-uncertainty cohort exhibit $+5.6025$ marks higher grading error than low-uncertainty scripts (14.8333 vs. 9.2308 marks, a 60.7\% error amplification), confirming reliable risk discrimination without ground-truth leakage.

\subsection{Selective Review & Workload Reduction}
Following the selective classification framework of Geifman and El-Yaniv~\cite{geifman2017selective}, we evaluate human-in-the-loop triage across 12 review budgets. Scripts exceeding an uncertainty threshold are escalated to human examiners (who award gold marks), while the remainder are graded autonomously.

\begin{table}[t]
\centering
\caption{Calibrated Selective Human Review Curves ($N = 37$ Booklets)}
\label{tab:selective_curves}
\begin{tabular}{cccccc}
\toprule
\textbf{Budget} & \textbf{Human} & \textbf{Autonomous} & \textbf{Saved} & \textbf{Residual MAE} & \textbf{Residual QWK} \\
\midrule
0.0\% & 0 & 37 & 100.0\% & 12.8649 & -0.0045 \\
10.0\% & 4 & 33 & 89.2\% & 11.4595 & -0.0043 \\
20.0\% & 7 & 30 & \textbf{81.1\%} & \textbf{10.3243} & 0.0224 \\
50.0\% & 18 & 19 & \textbf{51.4\%} & \textbf{5.0541} & 0.5887 \\
75.0\% & 28 & 9 & \textbf{24.3\%} & \textbf{1.8378} & \textbf{0.8339} \\
100.0\% & 37 & 0 & 0.0\% & 0.0000 & 1.0000 \\
\bottomrule
\end{tabular}
\end{table}

Table~\ref{tab:selective_curves} details the triage progression:
\begin{itemize}
    \item At a 20\% review budget (escalating only 7 scripts), residual MAE drops from 12.8649 to 10.3243, reducing manual grading workload by \textbf{81.1\%}.
    \item At a 50\% review budget, MAE drops to 5.0541, exceeding Oracle grading parity while cutting human grading effort in half (\textbf{51.4\% workload saved}).
    \item At a 75\% review budget, residual MAE drops to 1.8378, achieving substantial inter-rater reliability with human examiners ($\kappa = \textbf{0.8339}$).
\end{itemize}

\section{Discussion \& Threat Model}

\subsection{Operational Limitations & Platform Bounds}
During large-scale benchmarking, PaddleOCR encountered severe execution failures on the Windows platform due to a low-level oneDNN PIR C++ array attribute exception. We documented this failure transparently (6.39\% success rate, 93.61\% failure rate) rather than omitting or simulating results. This highlights the importance of hardware-software co-validation in document intelligence systems.

\subsection{Threats to Validity}
While our corpus contains 37 complete multi-page examination booklets (195 scanned pages, 407 subquestions), all exams were drawn from a single university institution and academic domain (Computer Science / Large Language Models). Handwriting variability, examination layout conventions, and mathematical symbology may differ in other STEM disciplines. Future work should expand this protocol to multi-institutional cohorts.

\section{Conclusion}
This paper presented the first comprehensive empirical investigation into assessment-critical handwriting recognition and error propagation in automated examination grading. By establishing a double-blind ground-truth protocol and executing a controlled $2 \times 2$ factorial experiment, we demonstrated that perceptual handwriting recognition is the primary bottleneck in automated grading ($F = 260.84, \eta_p^2 = 0.6443$), discovered the Preprocessing Paradox in Vision Transformers, resolved scan rotation fragility via hybrid layout continuation, and proved that input-side calibrated selective review can reduce human grading workload by over 81\% while maintaining assessment validity.

\section*{Acknowledgment}
The authors thank the Department of Computer Science and Engineering at IIIT Dharwad for supporting this institutional research.

\bibliographystyle{IEEEtran}
\bibliography{references}

\end{document}
